\section{Method}
\subsection{Data and cohort construction}
Source A is EBBU-2017 (\res{nAn} URLs: \res{nAbenign} legitimate, \res{nAphishing} phishing; repository \texttt{ebubekirbbr/pdd} at commit \texttt{90b394d}); the class is given by the file name. Source B is HP (\res{nBn} URLs), rebuilt by concatenating the \texttt{hp\_train/valid/test} files of \texttt{fariza25/PhishingURLDetection} at commit \texttt{ff9f1f5}; the original split is kept only as provenance. Because repository B has no codebook, we verified its label convention on its own copy of EBBU: joining it with source A on the raw URL gave \res{labelJoin} matches with \res{labelAgree}\% agreement, confirming 1 = phishing. All file hashes, sizes and observed counts are stored in a manifest; observed counts equal the expected ones.

URLs are canonicalised offline (no network access is possible during extraction): outer whitespace is trimmed; strings with control characters, schemes other than HTTP(S), invalid ports or unparsable hosts are quarantined (\res{quarA} in A, \res{quarB} in B); hosts are lower-cased and IDNA-encoded, while path, query and fragment keep their case and percent-encoding. Exact duplicates (same canonical key and label) are removed within each source (\res{dupA} in A, \res{dupB} in B), and keys with conflicting labels would be quarantined in both sources (none occurred). Each URL is grouped by its registered domain under the ICANN section of a frozen Public Suffix List snapshot; private suffixes are deliberately ignored so that, e.g., all \texttt{blogspot.com} sites form one conservative group. The two sources shared \res{sharedGroups} registered domains (only \res{exactOverlap} exact URLs); all URLs of these domains were removed from \emph{both} sources (\res{purgeA} from A, \res{purgeB} from B). The purge consults target identities only for this audit, never for learning, and costs mostly benign URLs (\res{lossBenignA}\% of A's and \res{lossBenignB}\% of B's benign URLs, versus \res{lossPhishA}\% and \res{lossPhishB}\% of phishing), because popular legitimate platforms appear in both. Table~\ref{tab:attrition} gives the full attrition. The resulting cohorts have \res{cAn} (A) and \res{cBn} (B) URLs over \res{groupsA} and \res{groupsB} domains; the largest group is \texttt{\res{topGroupA}} (\res{topGroupShareA}\%) in A and \texttt{\res{topGroupB}} (\res{topGroupShareB}\%) in B.

\begin{table}[t]
\caption{Cohort attrition (each URL counted once, at the first stage that removed it).}
\label{tab:attrition}
\centering\scriptsize
\begin{tabular}{llrrr}
\toprule
Source & Stage & URLs & Benign & Phishing \\
\midrule
EBBU-2017 & Loaded & 73,575 & 36,400 & 37,175 \\
EBBU-2017 & Parsed & 73,572 & 36,400 & 37,172 \\
EBBU-2017 & Deduplicated & 54,761 & 29,910 & 24,851 \\
EBBU-2017 & Conflicts removed & 54,761 & 29,910 & 24,851 \\
EBBU-2017 & Cohort (shared domains purged) & 46,769 & 22,167 & 24,602 \\
\midrule
HP & Loaded & 92,500 & 46,250 & 46,250 \\
HP & Parsed & 92,483 & 46,233 & 46,250 \\
HP & Deduplicated & 92,482 & 46,233 & 46,249 \\
HP & Conflicts removed & 92,482 & 46,233 & 46,249 \\
HP & Cohort (shared domains purged) & 75,077 & 29,758 & 45,319 \\
\bottomrule
\end{tabular}

\end{table}

\subsection{Splits}
Each source is split into FIT/VAL/TEST ($\approx$60/20/20\% of URLs) by a deterministic greedy stratified \emph{group} assignment (seed 20261115): whole domains are placed, largest first, into the partition whose per-class quota is most unfilled. Partitions share no domain, within or across sources, and keep the class prevalence (Table~\ref{tab:splits}). FIT is used for selection and fitting, VAL only for thresholds, and TEST once, after freezing.

\begin{table}[t]
\caption{Group-disjoint partitions.}
\label{tab:splits}
\centering\scriptsize
\begin{tabular}{llrrrr}
\toprule
Source & Part & URLs & Domains & \% URLs & \% phish \\
\midrule
EBBU & FIT & 28,061 & 8,018 & 60.0 & 52.6 \\
EBBU & VAL & 9,355 & 6,448 & 20.0 & 52.6 \\
EBBU & TEST & 9,353 & 6,447 & 20.0 & 52.6 \\
HP & FIT & 45,050 & 2,614 & 60.0 & 60.4 \\
HP & VAL & 15,014 & 2,060 & 20.0 & 60.4 \\
HP & TEST & 15,013 & 2,061 & 20.0 & 60.4 \\
\bottomrule
\end{tabular}

\end{table}

\subsection{Features}
We use 37 deterministic lexical features (lengths of URL, host, path, query and fragment; counts of 10 delimiter characters, digits, letters and non-ASCII code points; digit ratios; character entropy; longest alphanumeric run; host label and subdomain counts; IP host, explicit port and port number; user-info, query, HTTPS and missing-scheme flags; punycode labels; percent-encoded triplets; query pairs; path segments; and token counts of a fixed 10-word keyword list in the URL and in the host). Raw strings, domains, dataset identity and split labels never enter a model; no vocabulary or target encoding is learned. Each feature is a separate function, so a subset computes only what it needs, while parsing is always paid.

\subsection{Source-only feature selection}
All selectors use the same RF proxy (200 trees, depth 16, min.\ leaf 2, $\sqrt{p}$ features). FIT is divided into three stratified \emph{domain} folds; in each fold the proxy is trained on two folds and the permutation importance (PI) of every feature is measured on the third as the drop in average precision over five permutations. MEAN-PI ranks features by mean PI over the folds. STABLE-PI first counts, for each feature, the folds in which it is in the top $q{=}12$ \emph{and} has PI${>}0$ (votes), then ranks by votes, median PI, the median absolute deviation of its fold ranks, and feature ID. The budget $k\in\{8,16,24\}$ is chosen by nested cross-validation: in three outer domain folds of FIT, PI is recomputed from the outer-training part only, each classifier is trained on the top-$k$ features and scored on the outer hold-out; the smallest $k$ within 0.01 AP of the best mean is kept, separately for every source, classifier, selector and seed. PI is then recomputed on all of FIT for the final subset. To isolate the effect of stability from that of the budget, every STABLE-PI model is paired with a MEAN-PI model at the \emph{same} $k$ (matched-$k$ contrast).

\subsection{Classifiers, thresholds and evaluation}
RF uses 300 trees, depth 20, min.\ leaf 2, $\sqrt{p}$ features; LightGBM uses 300 rounds, learning rate 0.05, 31 leaves, no subsampling, $\lambda_2{=}1$ and deterministic CPU mode. Hyper-parameters are fixed a priori and are not claimed to be optimal~\cite{bergstra2012}. Models are trained on source FIT with seeds 17, 29 and 43. On source VAL we lock (i) the threshold maximising F1 (primary), (ii) the lowest threshold whose VAL false-positive rate is $\le$1\% and (iii) 0.5 (sensitivity). The same model and threshold are then applied to the source TEST (internal) and to the other source's TEST (cross), with no retraining or recalibration on the target. The transfer drop is F1(source TEST)$-$F1(target TEST).

The pre-specified primary outcome is target-TEST F1 of STABLE-PI minus MEAN-PI at matched $k$, for each classifier and direction (four contrasts). Uncertainty is a paired domain-cluster bootstrap (\res{nBoot} replicates): target domains are resampled with replacement, both arms and all three seeds use the same resample, and the per-replicate delta is averaged over seeds. These intervals are conditional on the fixed split and seeds, and we make no family-wise significance claim; an absolute F1 difference of 0.01 was fixed in advance as the smallest difference worth discussing.

\subsection{Cost and attributions}
Latency is measured on the seed-17 model of every configuration with the same interleaved VAL URLs of both sources: 10 warm-up and 300 (batch 1) or 30 (batch 1024) timed repetitions, reporting medians for parsing+feature extraction, prediction alone and end-to-end, with one and with four inference threads. SHAP values are computed with TreeExplainer on the actual classifiers (interventional, probability output, 100 source-FIT background rows) for 500 label-stratified source-TEST and 500 target-TEST URLs, for the seed-17 ALL and STABLE-PI models; additivity against the predicted phishing probability holds to \res{shapAddErr}.

\subsection{Protocol control}
Code, configurations, data splits, feature matrices and all fitted models of both directions were hashed into a freeze manifest (hash \texttt{\res{freezeHash}}) before the evaluator read any TEST label; the development loader cannot load TEST at all. Every number in this paper is generated from the saved predictions and re-computed from them as a check. This is a procedural control rather than a cryptographic guarantee.
